\section{The Canonical Scene}

\subsection{John slows the race at the doorway}

Mary Magdalene comes to the tomb while it is still dark, sees the stone
removed, and reports an inference: ``They have taken away the Lord out of the
sepulchre, and we know not where they have laid him'' (Jn 20:2). Her report
distinguishes the first observation from its first explanation. The body is
absent; removal is inferred; the agents and destination are unknown.

Peter and the other disciple run. The other arrives first, bends down, and sees
the linen cloths lying, but does not enter. Peter arrives, enters, and regards
both the linen cloths and the face-cloth. The other disciple then enters,
``saw, and believed.'' John immediately qualifies that belief: ``For as yet
they knew not the scripture, that he must rise again from the dead'' (20:9).

The progression is not a contest in spiritual rank. It is an exact narrative
sequence. One disciple first sees from outside; Peter crosses the threshold and
examines; the first arrival follows and believes. John's syntax makes the last
subject the subject of believing. Peter's role is indispensable to the scene,
but the narrator does not transfer the later verb backward.

\subsection{Luke preserves Peter's wonder}

Luke 24:12 gives a parallel in compressed form. Peter rises, runs, bends toward
the tomb, sees the linen cloths alone, and departs wondering at what happened.
Luke's \emph{thaumaz\=o} can accompany amazement, admiration, perplexity, or a
response not yet settled into understanding. Within Luke 24, the women report
the angelic message; the apostolic group initially treats it as idle talk;
Peter investigates; the Emmaus disciples later say that some went to the tomb
and found matters as the women said but did not see Jesus.

That sequence makes wonder an honest midpoint. Peter has evidence against
Mary's unexamined assumption and evidence that demands an explanation. Luke
does not say that he disbelieved the Resurrection; neither does Luke say that
he had already concluded it from the cloths.

\subsection{The other canonical accounts}

Matthew reports the stone, angel, empty place, guard, and later theft story.
Mark reports the moved stone, the young man, the place where Jesus had lain,
and the message to the disciples and Peter. Neither gives the cloth inventory.
Their silence neither contradicts John nor independently confirms the
arrangement. The theft polemic in Matthew supplies a live ancient question to
which later interpreters could direct John's details.

The Fourth Gospel itself supplies the closest control in Lazarus. Lazarus comes
from the tomb still bound in burial cloths, his face wrapped with a
\emph{soudarion}; Jesus orders others to unbind him (Jn 11:44). In chapter 20
the risen Jesus is absent and the cloths remain. The contrast belongs to John's
own narrative world: resuscitation returns Lazarus bound; resurrection does not
present Jesus as a corpse needing release.

